\chapter{Real-World Engineering Applications: OOP, File I/O, and Exceptions}
\label{ch:real_world_applications}

\section{Object-Oriented Programming (OOP) in Python}

Object-Oriented Programming models real-world hardware systems as computational objects combining internal state (attributes) with behaviors (methods).

\subsection{Classes, Instances, and the \texttt{self} Parameter}
In Python, a class is a blueprint for instantiating objects. The special method \texttt{\_\_init\_\_} serves as the object constructor. The explicit \texttt{self} parameter refers to the current instance invoking the method.

\begin{lstlisting}[caption={Modeling an Embedded Analog Sensor Class}]
class SensorNode:
    # Class-level attribute shared across all instances
    BUS_STANDARD = "I2C"
    
    def __init__(self, node_id: str, address: int, v_ref: float = 3.3):
        # Instance attributes
        self.node_id = node_id
        self.address = address
        self.v_ref = v_ref
        self._calibration_offset = 0.0  # Protected attribute by convention
        
    def read_raw_adc(self) -> int:
        """Simulate reading 10-bit ADC integer."""
        import random
        return random.randint(0, 1023)
        
    def get_voltage(self) -> float:
        raw = self.read_raw_adc()
        return (raw / 1023.0) * self.v_ref + self._calibration_offset
        
    def set_offset(self, offset: float):
        if abs(offset) > 0.5:
            raise ValueError("Calibration offset exceeds safe +/- 0.5V limit.")
        self._calibration_offset = offset
\end{lstlisting}

\subsection{Inheritance and Polymorphism}
Subclasses inherit attributes and methods from base classes and can override behavior using the \texttt{super()} proxy object:

\begin{lstlisting}[caption={Specialized Temperature Sensor Subclass}]
class TemperatureSensor(SensorNode):
    def __init__(self, node_id: str, address: int, sensitivity_mv_c: float = 10.0):
        super().__init__(node_id, address, v_ref=3.3)
        self.sensitivity = sensitivity_mv_c / 1000.0  # Convert to V/deg C
        
    def get_temperature_celsius(self) -> float:
        voltage = self.get_voltage()
        # LM35 characteristic: 10 mV / deg C
        return voltage / self.sensitivity
\end{lstlisting}

\section{Defensive File Handling and Context Managers}

\subsection{The \texttt{with} Statement (Context Manager Protocol)}
Engineering scripts must guarantee that operating system file descriptors and hardware communication sockets are promptly closed, even if runtime errors occur. The \texttt{with} statement enforces the Context Manager protocol (\texttt{\_\_enter\_\_} and \texttt{\_\_exit\_\_}).

\begin{lstlisting}[caption={Reading and Writing Telemetry Logs with CSV}]
import csv

# Writing structured engineering log data
log_file = "sensor_telemetry.csv"
records = [
    {"timestamp": "12:00:01", "node": "TMP_01", "temp": 24.5},
    {"timestamp": "12:00:02", "node": "TMP_01", "temp": 24.8},
    {"timestamp": "12:00:03", "node": "TMP_01", "temp": 25.1}
]

with open(log_file, mode="w", newline="", encoding="utf-8") as f:
    writer = csv.DictWriter(f, fieldnames=["timestamp", "node", "temp"])
    writer.writeheader()
    writer.writerows(records)

# Safe reading
with open(log_file, mode="r", encoding="utf-8") as f:
    reader = csv.DictReader(f)
    for row in reader:
        print(f"Time: {row['timestamp']} | Temp: {float(row['temp']):.1f} C")
\end{lstlisting}

\section{Robust Exception Handling}

Exceptions provide a structured mechanism for handling runtime errors without crashing mission-critical processes.

\begin{definition}[Exception Architecture]
The \texttt{try-except-else-finally} suite executes instructions with deterministic cleanup semantics:
\begin{itemize}[leftmargin=1.8em, itemsep=2pt]
    \item \textbf{\texttt{try}:} Encloses risky code that may raise an exception.
    \item \textbf{\texttt{except ExceptionType as err}:} Catches and handles specific exception instances.
    \item \textbf{\texttt{else}:} Executes if and only if no exception was raised in the \texttt{try} block.
    \item \textbf{\texttt{finally}:} Executes unconditionally under all circumstances (even after \texttt{return} or uncaught exceptions), ideal for releasing hardware locks.
\end{itemize}
\end{definition}

\begin{lstlisting}[caption={Comprehensive Exception Handling Pattern}]
def parse_telemetry_packet(raw_string: str) -> dict:
    try:
        parts = raw_string.strip().split(',')
        if len(parts) < 3:
            raise IndexError("Packet contains fewer fields than expected format.")
        
        node_id = parts[0]
        voltage = float(parts[1])
        status_code = int(parts[2])
        
    except ValueError as ve:
        print(f"Data corruption error: numeric conversion failed: {ve}")
        return {}
    except IndexError as ie:
        print(f"Malformed framing error: {ie}")
        return {}
    else:
        print("Packet parsed successfully with zero errors.")
        return {"node": node_id, "voltage": voltage, "status": status_code}
    finally:
        # Guarantee watchdog timer reset
        pass
\end{lstlisting}

\section{Real-World Case Study: Automated Diagnostic Log Processing}
Modern electronics systems deploy autonomous Python scripts to parse daily hardware crash dumps, calculate mean-time-between-failures (MTBF), generate analytical figures with Matplotlib, and dispatch diagnostic summaries to engineering teams.

\section{Chapter Summary and Review Questions}
\begin{enumerate}[leftmargin=1.8em, itemsep=4pt]
    \item Explain how Python's \texttt{with open(...) as f:} guarantees file closure even if an unhandled exception occurs inside the block.
    \item Contrast class attributes versus instance attributes. In what scenarios is a class attribute appropriate?
    \item What is the purpose of the \texttt{else} and \texttt{finally} clauses in a \texttt{try-except} construct?
    \item Design an object-oriented Python hierarchy representing electronic lab instruments (\texttt{PowerSupply}, \texttt{Oscilloscope}, \texttt{Multimeter}) derived from an abstract base class \texttt{LabInstrument}.
\end{enumerate}
